\makeabstract
    {
    Water-Cooling Acousto-Optics for a Paul Ion Trap
    }
    {
    Adrita Risha\textsuperscript{1}, Tiffany Payne-Igharoro\textsuperscript{1}, David Hanneke\textsuperscript{1}
    }
    {
    \textsuperscript{1} Department of Physics \& Astronomy, Amherst College, Amherst, MA
    }
    {
    Our lab uses a linear Paul trap to hold atomic and molecular ions for precision spectroscopy. Laser-cooled calcium ions are used to cool and detect trapped ions, requiring precise control of the frequencies of several laser beams. Acousto-optic modulators (AOMs) drive radio-frequency (RF) signals into crystals, producing acoustic waves that can diffract and manipulate the frequency of laser light. However, the RF power dissipated in an AOM generates heat, which can produce temperature gradients and affect the stability of the diffracted beam. So, we constructed a water-cooling system to thermally stabilize the AOMs in our laser system.
    }
    {
    We adapted a commercial computer water-cooling system for the AOMs. The system consisted of a water reservoir, pump, radiator and fans connected through flexible tubing. Two water-cooling blocks would sandwich each AOM. The two AOMs, one resonant and one detuned, were connected in series in one water loop. We monitored the AOM temperatures and subsequently characterized the optical system by measuring laser power, photodiode voltage, diffraction efficiency, and beam alignment. 
    }
    {
    Without active cooling, the AOM temperature increased to approximately 40 C during operation. With water cooling, the measured AOM temperature remained at room temperature, indicating that the system effectively removed any heat we could measure externally. After more realignment, the resonant AOM produced first-pass and second-pass diffraction efficiencies of approximately 74\% and 62\%, respectively. Although the external AOM temperatures were successfully stabilized, the measurements did not directly determine the temperature distribution within the AOM crystals. 
    }
    {
    The water-cooling system successfully prevented the large temperature increase previously observed on the AOMs and maintained both AOMs near room temperature during operation. However, because the cooling blocks contact the exterior of the AOM rather than the crystal directly, further measurements are required to determine whether the system also suppresses the internal temperature gradients.
    }

    \newpage
    